\documentclass{beamer}

% Кодировки (ВАЖНО для русского!)
\usepackage[T2A]{fontenc}
\usepackage[utf8]{inputenc}
\usepackage[russian]{babel}

\usepackage{graphicx}

% ---------------- СТРОГАЯ ТЕМА ----------------
\usetheme{Madrid}
\usecolortheme{default}

% Цветовая схема
\definecolor{darkblue}{RGB}{20,40,90}
\setbeamercolor{title}{fg=white}
\setbeamercolor{subtitle}{fg=white}
\setbeamercolor{block title}{bg=darkblue!15, fg=black}
\setbeamercolor{block body}{bg=gray!5, fg=black}

% Убираем навигацию
\setbeamertemplate{navigation symbols}{}
\setbeamertemplate{footline}[frame number]

% ---------------- ИНФОРМАЦИЯ ----------------
\title{Реализация системы для автоматизации процесса введения новых сотрудников в работу организации}
\author{Мамедова Фирана Ильдыровна}
\institute{Иркутский государственный университет}
\date{4 мая 2026 г.}

\begin{document}

% ---------------- ТИТУЛЬНЫЙ СЛАЙД ----------------
\begin{frame}
    \titlepage
    
    \vspace{0.5cm}
    \begin{center}
        \begin{tabular}{rl}
            \textbf{Руководитель:} & Рябец Л. В. \\[2pt]
            \textbf{Консультант:} & Специалист отдела разработки веб-решений \\
                                   & ЭН+ ДИДЖИТАЛ ООО\\
                                   & Афанасьев Евгений Андреевич \\
        \end{tabular}
    \end{center}
\end{frame}

% ---------------- СЛАЙД 2: АКТУАЛЬНОСТЬ ----------------
\begin{frame}
    \frametitle{Актуальность}
    
    \textbf{«Боль» пользователя:}
    \begin{itemize}
        \item Новичок теряется в разрозненных инструкциях
        \item Наставник тратит до 40\% времени на повторяющиеся вопросы
        \item Нет прозрачного контроля прогресса
    \end{itemize}
    
    \vspace{0.4cm}
    \textbf{Решение:}
    
    Единое веб-приложение для структурированного онбординга с отслеживанием прогресса
\end{frame}

% ---------------- СЛАЙД 3: ЦЕЛЬ И ЗАДАЧИ ----------------
\begin{frame}
    \frametitle{Цель и задачи}
    
    \textbf{Цель:}
    
    Разработать веб-приложение для автоматизации адаптации новых сотрудников ЭН+ Диджитал
    
    \vspace{0.4cm}
    \textbf{Задачи:}
    \begin{enumerate}
        \item Анализ предметной области и существующих решений
        \item Формирование требований к системе
        \item Разработка технического задания
        \item Проектирование архитектуры и структуры данных
        \item Реализация клиентской и серверной частей
        \item Тестирование и оценка соответствия требованиям
        \item Подготовка системы к внедрению
        \item Анализ результатов и направления развития
    \end{enumerate}
\end{frame}

% ---------------- СЛАЙД 4: АНАЛИЗ И ТРЕБОВАНИЯ ----------------
\begin{frame}
    \frametitle{Анализ и требования}
    
    \textbf{Что сделано:}
    \begin{itemize}
        \item Сравнительный анализ 5 коммерческих LMS/HR-систем
        \item Выявлена нецелесообразность готовых решений для ЭН+ Диджитал
        \item Сформированы 20+ функциональных и нефункциональных требований
        \item Разработано ТЗ с ролевой моделью и сценариями
    \end{itemize}
    
    \vspace{0.3cm}
    \textbf{Личный вклад:}
    
    Самостоятельно проведён анализ, подготовлены таблицы сравнения, сформулированы критерии применимости
\end{frame}

% ---------------- СЛАЙД 5: ПРОЕКТИРОВАНИЕ ----------------
\begin{frame}
    \frametitle{Проектирование системы}
    
    \textbf{Архитектура:}
    \begin{itemize}
        \item N-tier: Vue.js 3 ↔ ASP.NET Core 9 ↔ MS SQL Server
        \item JWT-аутентификация + BCrypt для хеширования
        \item 19 сущностей БД, гибкая ролевая модель
    \end{itemize}
    
    \vspace{0.3cm}
    \textbf{Особенности:}
    \begin{itemize}
        \item Адаптивная модель расчёта прогресса и геймификации
        \item Интеграционные контракты с RIMS, ПСО, SSO
        \item Документация API через Swagger (OpenAPI)
    \end{itemize}
\end{frame}

% ---------------- СЛАЙД 6: РЕАЛИЗАЦИЯ — BACKEND ----------------
\begin{frame}
    \frametitle{Реализация: серверная часть}
    
    \textbf{Основные сервисы (9 шт):}
    \begin{itemize}
        \item \texttt{GamificationService}: XP, уровни, достижения, рейтинг
        \item \texttt{AiMentorService}: Groq API LLaMA, контекст FAQ + PDF
        \item \texttt{ProgressAnalyticsService}: AI-аналитика, рекомендации
        \item \texttt{LearningPathService}: управление путём обучения
        \item \texttt{NotificationService}: уведомления (email + WebSocket)
        \item \texttt{EmailService}: отправка email-сообщений
        \item \texttt{ReportExportService}: экспорт PDF (QuestPDF) + Excel (ClosedXML)
        \item \texttt{RimsIntegrationService}: интеграция с RIMS API
        \item \texttt{KeycloakAdminService}: управление пользователями
    \end{itemize}
    
    \vspace{0.15cm}
    \textbf{Безопасность: JWT + BCrypt + роли + аудит логов}
    
    \vspace{0.15cm}
    \textbf{Инфраструктура: 18 контроллеров, Entity Framework Core, MS SQL Server}
\end{frame}

% ---------------- СЛАЙД 7: РЕАЛИЗАЦИЯ — FRONTEND ----------------
\begin{frame}
    \frametitle{Реализация: клиентская часть}
    
    \textbf{Компоненты (8 шт):}
    \begin{itemize}
        \item \texttt{ProgressMap.vue}: визуализация пути онбординга
        \item \texttt{TestView.vue}: пошаговое прохождение тестов с вопросами
        \item \texttt{AiChatAssistant.vue}: чат с AI-наставником (Groq API)
        \item \texttt{AnalyticsDashboard.vue}: интерактивный дашборд аналитики
        \item \texttt{Leaderboard.vue}: таблица лидеров с рейтингом
        \item \texttt{QuillEditor.vue}: редактор форматированного контента
        \item \texttt{LearningPath.vue}: интерактивная дорожная карта обучения
        \item \texttt{NotificationBell.vue}: real-time уведомления
    \end{itemize}
    
    \vspace{0.2cm}
    \textbf{Стек: Vue 3 + Vite + Tailwind CSS + SignalR (WebSocket)}
    
    \vspace{0.2cm}
    \textbf{Геймификация:}
    
    Визуальные эффекты (canvas-confetti) при завершении модулей, система XP, уровни, достижения, real-time рейтинг
\end{frame}

% ---------------- СЛАЙД 8: ТЕСТИРОВАНИЕ И ВНЕДРЕНИЕ (ИСПРАВЛЕННЫЙ) ----------------
\begin{frame}
    \frametitle{Тестирование и подготовка к внедрению}
    
    \textbf{Проведено:}
    \begin{itemize}
        \item \textbf{Unit-тесты backend:} 25--28 тестов (GamificationService, AuthenticationService, ModuleService)
        \item \textbf{Unit-тесты frontend:} 19 тестов валидаторов и состояния приложения
        \item \textbf{Integration-тесты:} 7--8 тестов сценариев аутентификации
        \item Проверка API через Swagger: 18 контроллеров
        \item Кросс-браузерная проверка (Chrome, Firefox, Edge)
        \item Адаптивность: мобильные, планшеты, десктоп
    \end{itemize}
    
    \vspace{0.3cm}
    \textbf{Готовность к развёртыванию:}
    \begin{itemize}
        \item Docker-образы для frontend (nginx) и backend
        \item Автоматические миграции БД при старте
        \item Инструкции для администратора и пользователя
    \end{itemize}
\end{frame}

% ---------------- СЛАЙД 9: РЕЗУЛЬТАТЫ И ПЕРСПЕКТИВЫ ----------------
\begin{frame}
    \frametitle{Результаты и направления развития}
    
    \textbf{Достигнуто:}
    \begin{itemize}
        \item Полнофункциональная система, готовая к пилотному внедрению
        \item Сокращение ручного труда наставников
        \item Положительная обратная связь от тестировщиков
    \end{itemize}
    
    \vspace{0.3cm}
    \textbf{Перспективы:}
    \begin{itemize}
        \item Полная интеграция с корпоративными системами
        \item Расширение AI-аналитики
        \item Мобильное приложение для обучения 
    \end{itemize}
\end{frame}

% ---------------- СЛАЙД 10: ВИДЕО-ДЕМОНСТРАЦИЯ ----------------
\begin{frame}
    \frametitle{Демонстрация системы}
    
    \centering
    \large \textbf{Видео: ключевой сценарий использования системы} \\[0.5em]
    \small \textit{(2–3 минуты)} \\[1em]
    
\end{frame}

% ---------------- СЛАЙД 11: ЗАКЛЮЧЕНИЕ ----------------
\begin{frame}
    \frametitle{Заключение}
    
    \textbf{Выполненные задачи:}
    \begin{enumerate}
        \item Проведён анализ предметной области и существующих решений
        \item Сформированы функциональные и нефункциональные требования
        \item Разработано техническое задание на создание системы
        \item Спроектирована архитектура приложения и структура данных
        \item Реализованы клиентская и серверная части программного продукта
        \item Проведено тестирование и оценка соответствия требованиям
        \item Система подготовлена для внедрения в организации
        \item Проведён анализ результатов и определены направления развития
    \end{enumerate}
\end{frame}

% ---------------- СЛАЙД 12: СТРУКТУРА ВКР ----------------
\begin{frame}
    \frametitle{Структура выпускной квалификационной работы}
    
    \centering
    \large
    \textit{[Сюда вставить скриншот оглавления ВКР с номерами страниц]}
    
    \vspace{0.5cm}
    \small
    \textbf{Примечание:}
    
    На итоговой защите этот слайд необходимо убрать
\end{frame}

\end{document}
